\begingroup

\section{The SGA trace-period source: exact coefficient and source corrections}


\subsection{The SGA trace-period source: exact coefficient and source corrections}
\label{sec:sga-trace-period-source-corrections}

The complete adjoining source is \emph{SGA trace, the retained diagonal,
and the arithmetic period determinant}, with its original equations
1--62, including 21a, 28a and 54a--54c.  Its delivered Markdown and
generated LaTeX remain byte-exact source witnesses.  Six explicitly
recorded replacements produce the corrected full LaTeX.  The changes
retain every proof paragraph and all original polynomial coefficients,
operators, source metrics, ray phases and orientations.  This introduction
states the coefficient types and the source qualifications needed when
reading the unaltered witness.

\subsubsection{The relative trace and its coefficient specialization}

For a commutative ring $R$ and the original monic polynomial $\chi$ of
degree $q>0$, set
\[
 T=R[t],\qquad H(S)=\chi(S)-t,\qquad E_t=T[S]/(H).
 \tag{SGAC1}
\]
The ordered basis $1,S,\ldots,S^{q-1}$ makes $E_t$ a free
$T$-module of rank $q$.  Division by the monic $H$ proves existence
of the displayed representative; a nonzero multiple of $H$ cannot
have degree below $q$, proving uniqueness.  Thus multiplication by
$P$ is a finite-free $T$-endomorphism, whose trace has the exact type
\[
 \operatorname{Tr}_{E_t/R[t]}(M_P)\in R[t].
 \tag{SGAC2}
\]
In the delivered equation (1), the base $R$ must therefore be read as
$R[t]$, already printed correctly in equation (21).  Finiteness over
$R[t]$ alone does not assert finiteness over $R$.

Retain the source's functional
$\lambda_t([P])=[S^{q-1}]\operatorname{rem}_H P$ and its polynomial
divided difference $\mathcal C_H(X,Y)=(H(X)-H(Y))/(X-Y)$ in
$E_t\otimes_T E_t$.  The displayed fraction is polynomial division.
Its $X^{q-1}$ coefficient is one and
$(X-Y)\mathcal C_H=0$.  Since $P(X)-P(Y)$ is divisible by $X-Y$,
these two identities give
\[
 (\lambda_t\otimes1)((P\otimes1)\mathcal C_H)=P,
 \qquad \mu(\mathcal C_H)=\chi'(S).
 \tag{SGAC3}
\]
Consequently $\mathcal C_H$ is the actual coevaluation tensor for
the residue pairing $(x,y)\mapsto\lambda_t(xy)$.  Its contraction
with multiplication by $P$ is
\[
 T\xrightarrow{1\mapsto\mathcal C_H}E_t\otimes_T E_t
 \xrightarrow{M_P\otimes1}E_t\otimes_T E_t
 \xrightarrow{\mu}E_t\xrightarrow{\lambda_t}T,
 \qquad
 \operatorname{Tr}_{E_t/T}M_P=\lambda_t(\chi'P).
 \tag{SGAC4}
\]
No derivative or discriminant is inverted by this composite.

For the parameter-universal polynomial
$\chi_{\mathbf c}(S)=S^q+\sum_{a=0}^{q-1}c_aS^a$, the trace bases
in delivered equations (2) and (35) are
$\mathbb C[\mathbf c,t]$.  In particular the complete formula is
\[
 \mathcal F(\mathbf c,t)
 =\operatorname{Tr}_{E_t/\mathbb C[\mathbf c,t]}(M_{\Phi_t}),
 \quad
 \Phi_t(S)=\frac{S^{q+1}}{q+1}
       +\sum_{a=0}^{q-1}\frac{c_aS^{a+1}}{a+1}-tS.
 \tag{SGAC5}
\]
For any coefficient evaluation $f:T\to T'$, monic division gives
the explicit basis-preserving isomorphism
$E_t\otimes_TT'\to T'[S]/(f(H))$ by
$[P]\otimes a\mapsto a[f(P)]$.  It has the inverse that sends
each of the same $q$ basis elements to its tensor basis element.
Every multiplication-matrix entry evaluates under $f$; summing its
diagonal therefore proves
$f(\operatorname{Tr}_{E_t/T}M_P)
=\operatorname{Tr}_{E_t\otimes_TT'/T'}M_{f(P)}$.
This is the exact family-to-fibre trace map.  At specified complex
coefficient and parameter values it gives the finite-dimensional
$\mathbb C$-linear trace, retaining the marked original polynomial.

\subsubsection{The differential coefficient ring and the metric paragraph}

In both terms of delivered equation (28), the universal ring is
$\mathbb C[\mathbf c,u,t,S]$.  The differential is the same
coefficient-linear map
\[
 D(P)=\bigl(uP'+(\chi_{\mathbf c}-t)P\bigr)dS,
 \qquad
 [\mathbb C[\mathbf c,u,t,S]\xrightarrow{D}
      \mathbb C[\mathbf c,u,t,S]dS].
 \tag{SGAC6}
\]
At the fixed original coefficient value this specializes to the
source's $\mathbb C[u,t,S]$ complex.  Differentiation in $S$ and
coefficient evaluation commute, so this is an actual map of the
two-term complexes.  The leading monic term gives the finite
reduction: for a term of degree at least $q$, subtract the corresponding
coefficient multiple of $D(S^m)$; its derivative contribution has
degree $m-1$, while its leading polynomial term has degree $m+q$.
Each subtraction lowers the unreduced leading degree and uses no
division by a parameter.  The resulting remainder is unique because
the leading term of $D(P)$ for nonzero $P$ has degree $\deg P+q$.
Every subtraction commutes with coefficient evaluation.  The range
of $D$ retains its source type as a linear relation; for example
$D(S)=(u+S(\chi_{\mathbf c}-t))dS$ gives the original relation
$[S(\chi_{\mathbf c}-t)]=-u[1]$.

The generated paragraph following equation (50) suffered Markdown
emphasis damage.  Its exact mathematical sentence is:
\begin{quote}
Conversely $\Pi^*\Pi$ is the Gram for the separately specified
coordinate metric $I_q$, and its comparison operator is
$G_N^{-1}\Pi^*\Pi$.
\end{quote}
Here $u\ne0$ and $\Pi=\Pi(\mathbf c,u,t)$ is the original period
matrix.  Its determinant is nonzero by the complete evaluation in
equations (38)--(41) of the adjoining note.  With the original positive
Gram $G_N$ retained, the metric transport is unchanged:
\[
 \widetilde G_N=\Pi^{-*}G_N\Pi^{-1},\qquad
 \Pi^*\widetilde G_N\Pi=G_N,
 \qquad G_N(G_N^{-1}\Pi^*\Pi)=\Pi^*\Pi.
 \tag{SGAC7}
\]
The two equalities are direct multiplication, with
$\Pi^{-*}=(\Pi^{-1})^*$.  Thus they specify both the original-metric
isometry and the comparison with the separately stated coordinate
metric, with all Gram factors retained.

\subsubsection{The independently inspected primary source}

The current independent primary witness is the original French
SGA~5 scan, printed pages 120--130 and 136--137.  The expected English
master was not located in the bounded search.  The delivered note's
statements about its author's supplied-master reading remain historical
source statements; this integration does not count the missing English
bytes as another complete reading.

Within SGA~5's noetherian setup, formula (6.8.4) assumes that the
ambient smooth map is proper.  Formula (6.8.5) uses the finite-flat
composite, regular immersion and smooth ambient map; ambient properness
is not an additional hypothesis of that formula.  Its printed
qualification concerning the proof/reference is retained in the
complete primary review.  The monic calculation SGAC1--SGAC4 supplies
the displayed trace identity over arbitrary commutative $R$ directly.

Printed page 128 has three $N_{Y/S}$ tokens where the preceding
definition and its Hom isomorphism give the conormal
$N_{Y/X}=I/I^2$.  This is an annotated source typography correction.
In the present monic model $I=(H)$ and the exact conormal differential is
\[
 I/I^2\longrightarrow E_t\,dS,\qquad
 [H P]\longmapsto[P]\chi'(S)dS.
 \tag{SGAC8}
\]
It is well-defined because changing $HP$ by $H^2Q$ changes its
derivative by a term divisible by $H$ after reduction.  Equivalently
$d(HP)=P\chi'dS+H\,dP$ reduces to the displayed image.  The
generator $[H]$ and the derivative $\chi'$ retain their actual
roles in the trace composite SGAC4.  Literal reference scans remain
reference-only; this correction does not silently alter them.

The six exact source replacements, their occurrence counts and
intermediate hashes are recorded in
\nolinkurl{cumulative_sga_embedding_corrections_public_20260913.json}.
The dictionary also binds the unchanged delivered TeX and its complete
corrected descendant.  It records no changed mathematical test result
and makes no additional arithmetic growth claim.


\endgroup
